\SourcePage{355}{360}
\section{Electron scattering in an external field}

Elastic scattering of an electron in a constant external field is the simplest process that already occurs in first-order perturbation theory (the first Born approximation). It is represented by a diagram with one vertex,
\begin{equation}
\begin{tikzpicture}[baseline=-.45ex,x=.72cm,y=.72cm,>=Stealth,line width=.65pt,font=\small]
\coordinate (v) at (0,0);
\draw[->] (1.15,-.82) node[below right] {$p$} -- (v);
\draw[->] (v) -- (-1.15,-.82) node[below left] {$p'$};
\draw[->] (0,1.05) -- (v) node[midway,right=2pt] {$k$};
\end{tikzpicture}
\tag{80.1}
\end{equation}
where $p$ and $p'$ are the initial and final electron four-momenta, and $q=p'-p$. Since the electron energy is conserved in scattering by a constant field ($\varepsilon=\varepsilon'$), $q=(0,\mathbf q)$\footnote{In the case of an external field, such a diagram is of course not forbidden by four-momentum conservation (as was diagram (73.19) with a real photon): unlike the square of the four-momentum of a real photon, $q^2$ need not vanish; the component with the required $q$ is selected automatically from the Fourier integral representing the external field.}.

The corresponding scattering amplitude is
\begin{equation}
M_{fi}=-e\bar u(p')[\gamma A^{(e)}(\mathbf q)]u(p),
\tag{80.2}
\end{equation}
where $A^{(e)}(\mathbf q)$ is the spatial Fourier component of the external field. According to (64.26), the scattering cross section is
\begin{equation}
d\sigma=\frac{1}{16\pi^2}|M_{fi}|^2\,do'.
\tag{80.3}
\end{equation}

For an electrostatic field $A^{(e)}=(A_0^{(e)},0)$, and hence
\begin{equation}
M_{fi}=-e\bar u(p')\gamma^0u(p)A_0^{(e)}(\mathbf q)
=-eu^*(p')u(p)A_0^{(e)}(\mathbf q).
\tag{80.4}
\end{equation}
In the nonrelativistic case, the plane-wave bispinor amplitudes $u(p)$ reduce to nonrelativistic (two-component) amplitudes. For scattering without a change of polarization, this is a quantity independent of $\mathbf p$; under our normalization convention, $u^*u=2m$. Taking this into account gives
\[
d\sigma=\left|\frac{m}{2\pi}U(\mathbf q)\right|^2do',
\]

\SourcePage{356}{361}
where $U(\mathbf q)=eA_0^{(e)}(\mathbf q)$ is the Fourier component of the electron's potential energy in the field; this expression agrees with the familiar Born formula (Vol.~III, (126.7)).

In the general relativistic case, the scattering cross section for unpolarized electrons is obtained by averaging $|M_{fi}|^2$ over the initial polarizations and summing over the final ones, that is, by forming
\[
\frac{1}{2}\sum_{\text{pol}}|M_{fi}|^2,
\]
where the sum is over the spin directions of the initial and final electrons; the factor $1/2$ converts one of these summations into an average. The rules set out in \S~65 give
\[
\begin{aligned}
\frac{1}{2}\sum_{\text{pol}}|M_{fi}|^2
&=2\Tr\rho'(\gamma A_0^{(e)})\rho(\gamma A_0^{(e)*})=\\
&=\frac{1}{2}|A_0^{(e)}(\mathbf q)|^2
\Tr(m+\gamma p')\gamma^0(m+\gamma p)\gamma^0.
\end{aligned}
\]
To evaluate the trace, note that $\gamma^0(\gamma p)\gamma^0=\gamma\widetilde p$, where $\widetilde p=(\varepsilon,-\mathbf p)$, and therefore
\[
\begin{aligned}
\frac{1}{4}\Tr(m+\gamma p')\gamma^0(m+\gamma p)\gamma^0
&=\frac{1}{4}\Tr(m+\gamma p')(m+\gamma\widetilde p)=\\
&=m^2+p'\widetilde p=\varepsilon^2+m^2+\mathbf p\mathbf p'
=2\varepsilon^2-\mathbf q^2/2.
\end{aligned}
\]
It follows that the cross section is
\begin{equation}
d\sigma=\frac{e^2|A_0^{(e)}(\mathbf q)|^2}{4\pi^2}\,
\varepsilon^2\left(1-\frac{\mathbf q^2}{4\varepsilon^2}\right)do'.
\tag{80.5}
\end{equation}

For the field produced by a static charge distribution of density $\rho(\mathbf r)$, we have
\begin{equation}
A_0^{(e)}(\mathbf q)=\frac{4\pi\rho(\mathbf q)}{\mathbf q^2},
\tag{80.6}
\end{equation}
where $\rho(\mathbf q)$ is the Fourier transform of the distribution $\rho(\mathbf r)$ (the form factor). In particular, for the Coulomb field of a point charge $Ze$, $\rho(\mathbf q)=Ze$. The scattering cross section is then
\begin{equation}
d\sigma=do'\frac{4(Ze^2)^2\varepsilon^2}{\mathbf q^4}
\left(1-\frac{\mathbf q^2}{4\varepsilon^2}\right)
\tag{80.7}
\end{equation}
(\emph{N.~F.~Mott}, 1929). The square
\[
\mathbf q^2=4\mathbf p^2\sin^2(\theta/2),
\]

\SourcePage{357}{362}
where $\theta$ is the scattering angle. The expression preceding the parentheses may therefore be called, from its angular dependence, the Rutherford cross section:
\begin{equation}
d\sigma_{\mathrm R}=do'\frac{4(Ze^2)^2\varepsilon^2}{\mathbf q^4}
=do'\frac{4(Ze^2)^2\varepsilon^2}{4\mathbf p^4}\sin^{-4}\frac{\theta}{2}
\tag{80.8}
\end{equation}
(in the nonrelativistic limit, the factor $\varepsilon^2/\mathbf p^4\to1/(m^2v^4)$). Thus\footnote{The difference between $d\sigma$ and $d\sigma_{\mathrm R}$ expressed by this formula is specific to particles of spin $1/2$. For the scattering of spin-0 particles (if their motion in an electromagnetic field were described by a wave equation), one would obtain $d\sigma=d\sigma_{\mathrm R}$. At first sight it may seem strange that the factor expressing this purely quantum effect contains no $\hbar$. It must be remembered, however, that the condition for applicability of the Born approximation ($e^2/(\hbar v)\ll1$) is opposite to the condition for quasiclassical motion in a Coulomb field; consequently, the classical limit cannot be taken in formula (80.9).},
\begin{equation}
d\sigma=d\sigma_{\mathrm R}\left(1-v^2\sin^2\frac{\theta}{2}\right).
\tag{80.9}
\end{equation}
Notice that in the ultrarelativistic case the angular distribution differs from the nonrelativistic one by a strong suppression of backward scattering (as $\theta\to\pi$: $d\sigma/d\sigma_{\mathrm R}\to m^2/\varepsilon^2$).

For small-angle scattering in the ultrarelativistic case, (80.7) gives
\begin{equation}
d\sigma=\frac{4(Ze^2)^2}{\varepsilon^4\theta^4}\,do'.
\tag{80.10}
\end{equation}
Although this formula was obtained in the Born approximation (that is, under the assumption $Ze^2\ll1$), it remains valid (for angles $\theta\lesssim m/\varepsilon$) also when $Ze^2\sim1$. This can be verified using the ultrarelativistic wave function $\psi_{\mathbf p}^{(+)}$ (39.10), which is exact in $Ze^2$. This solution, valid in the region (39.2), naturally remains valid in the asymptotic region at arbitrarily large $r$. There,
\[
F\propto1+\text{const}\cdot e^{i(pr-\mathbf p\mathbf r)},
\qquad
\frac{\alpha\nabla F}{\varepsilon}\sim1-\cos\theta\sim\theta^2\ll1,
\]
so the correction term remains small, as expected. The wave function of the form $e^{ipr}F$, being identical in form to the nonrelativistic function (with the evident change of parameters), has the same asymptotic form; the Rutherford expression therefore results for the cross section as well.

The cross section for scattering of electrons with arbitrary polarization could be calculated by the general rules using the density matrix (29.13). In the present case, however,

\SourcePage{358}{363}
the result can be obtained less laboriously by writing the bispinor amplitudes $u(p')$ and $u(p)$ in the form (23.9); multiplication gives
\[
u^*(p')u(p)=w'^*\{\varepsilon+m+(\varepsilon-m)
(\mathbf n'\boldsymbol\sigma)(\mathbf n\boldsymbol\sigma)\}w,
\]
or, using formula (33.5),
\begin{equation}
u^*(p')u(p)=w'^*\widehat f w,
\tag{80.11}
\end{equation}
where\footnote{The definitions of $\widehat f$ here and in \S~37, \S~38 differ by an overall factor.}
\begin{equation}
\begin{gathered}
\widehat f=A+B\boldsymbol\nu\boldsymbol\sigma,\\
A=(\varepsilon+m)+(\varepsilon-m)\cos\theta,
\qquad B=-i(\varepsilon-m)\sin\theta,\\
\boldsymbol\nu=\frac{\mathbf n\times\mathbf n'}{\sin\theta}.
\end{gathered}
\tag{80.12}
\end{equation}
The two-component quantity (three-spinor) $w$ is the nonrelativistic spin wave function of the electron. Partially polarized states are therefore introduced by replacing the products $w_\alpha w_\beta^*$ ($\alpha$, $\beta$ are spinor indices) by the nonrelativistic two-by-two density matrix $\rho_{\alpha\beta}$. Thus one must replace
\[
|M_{fi}|^2\to e^2|A_0^{(e)}(\mathbf q)|^2
\Tr\rho(A-B\boldsymbol\nu\boldsymbol\sigma)
\rho'(A+B\boldsymbol\nu\boldsymbol\sigma),
\]
where
\[
\rho=\frac{1}{2}(1+\boldsymbol\sigma\boldsymbol\zeta),
\qquad
\rho'=\frac{1}{2}(1+\boldsymbol\sigma\boldsymbol\zeta'),
\]
and $\boldsymbol\zeta$ and $\boldsymbol\zeta'$ are the initial polarization vector and the final polarization vector selected by the detector. Evaluation of the trace gives
\begin{equation}
d\sigma=d\sigma_0\left\{1+
\frac{(A^2-|B|^2)(\boldsymbol\zeta\boldsymbol\zeta')
+2|B|^2(\boldsymbol\nu\boldsymbol\zeta)(\boldsymbol\nu\boldsymbol\zeta')
+2A|B|\boldsymbol\nu(\boldsymbol\zeta\times\boldsymbol\zeta')}
{A^2+|B|^2}\right\},
\tag{80.13}
\end{equation}
where $d\sigma_0$ is the cross section for scattering of unpolarized electrons.

Writing the expression in braces in (80.13) as $\{1+\boldsymbol\zeta^{(f)}\boldsymbol\zeta'\}$, we find the polarization of the final electron itself (as distinct from the detected polarization $\boldsymbol\zeta'$; see \S~65)\footnote{Formula (80.14) corresponds to the formula found in Problem 1 (see Vol.~III, \S~140) and follows from it when $A$ is real and $B$ imaginary.}:
\begin{equation}
\boldsymbol\zeta^{(f)}
=\frac{(A^2-|B|^2)\boldsymbol\zeta
+2|B|^2(\boldsymbol\nu\boldsymbol\zeta)\boldsymbol\nu
+2A|B|(\boldsymbol\nu\times\boldsymbol\zeta)}
{A^2+|B|^2}.
\tag{80.14}
\end{equation}
We see that the scattered electrons are polarized only if the incident electrons are polarized. This is a general property of the first Born approximation (cf. Vol.~III, \S~140).

\SourcePage{359}{364}
In the nonrelativistic case ($\varepsilon\to m$), (80.14) gives $\boldsymbol\zeta^{(f)}=\boldsymbol\zeta$, that is, the electron retains its polarization upon scattering (a natural consequence of neglecting the spin--orbit interaction).

In the opposite, ultrarelativistic case, we have
\[
A=\varepsilon(1+\cos\theta),
\qquad
B=-i\varepsilon\sin\theta
\]
(in agreement with the general formula (38.2)).

If the incident electron then has a definite helicity ($\boldsymbol\zeta=2\lambda\mathbf n$, $\lambda=\pm1/2$), straightforward simplification of (80.14) gives
\[
\boldsymbol\zeta^{(f)}=2\lambda\mathbf n'.
\]
In other words, after scattering the electron remains in a helicity state and retains its previous helicity value $\lambda$.

As already explained in \S~38, this property is connected with the fact that, when the mass is neglected, the Dirac equation in the spinor representation separates into two independent equations for the functions $\xi$ and $\eta$. The result has a more general significance as well, because the current
\[
\widehat j=(\xi^*\xi+\eta^*\eta,\,
\xi^*\boldsymbol\sigma\xi-\eta^*\boldsymbol\sigma\eta),
\]
and hence the electromagnetic perturbation operator $\widehat V=e\widehat j\widehat A$, contain no terms mixing $\xi$ and $\eta$ and therefore have no matrix elements for transitions between $\xi$ and $\eta$ states. It follows that if an ultrarelativistic electron has a definite helicity (that is, either $\xi$ or $\eta$ is nonzero), then this helicity is conserved in interaction processes in the approximation in which the electron mass is neglected completely.
